\chapter{Events and Their Probabilities}

\section{Experiment, Sample Space, and Random Event}

\begin{definition}[Sample Space]\index{sample space}
  The set for all possible outcomes of an random experiment (if we could know all possible outcomes) is known as the \emph{sample space}.
\end{definition}

\begin{definition}[Random Event]\index{random event}\index{event}
  Any subset of the sample space \(S\) (or \(\Omega\)) is known as an \emph{random event} or \emph{event}.
\end{definition}
Events are usually denoted by capital letters.

\begin{itemize}
  \index{inevitable event}\index{impossible event}
  \item Inevitable event: \(S\) or \(\Omega\)
  \item Impossible event: \(\varnothing\)
\end{itemize}
\begin{remark}
An event with probability 0 is not necessarily an impossible event, and an event with probability 1 is not necessarily an inevitable event.
\end{remark}
\begin{example}
  Consider generating a random number \(X\) uniformly distributed on the interval \([0, 1]\). The event \(A = \{X = 0.5\}\) has probability 0, but it is not an impossible event since it can occur.
\end{example}

Some events can be expressed in terms of other events. For example, if \(A\) and \(B\) are two events, then we can define the following events:
\begin{enumerate}
  \index{subset}\index{union}\index{intersection}\index{set difference}\index{disjoint events}\index{complement}
  \item \(A \subset B\) means if event \(A\) occurs, then event \(B\) also occurs.
  \item \(A = B\) means \(A \subset B\) and \(B \subset A\), i.e., \(A\) and \(B\) are the same event.
  \item \(A \cup B\) occurs if either event \(A\) or event \(B\) (or both) occur.
  \item \(A B\) or \(A \cap B\) occurs if both event \(A\) and event \(B\) occur.
  \item \(A \setminus B\) occurs if event \(A\) occurs but event \(B\) does not occur.
  \item \(A B = \varnothing\) means events \(A\) and \(B\) are \emph{disjoint} or \emph{mutually exclusive}, i.e., they cannot occur at the same time.
  \item \(A^c\) or \(\overline{A}\) is the \emph{complement} of event \(A\), \(A \cup \overline{A} = \Omega\) and \(A \overline{A} = \varnothing\).
\end{enumerate}

\begin{example}
  Select points randomly from unit circle. Let
  \[
  A_n = \left\{(x,y) \vert x^2 + y^2 < \frac{1}{n^2}\right\}, n = 1, 2, \dots
  \]
  then
  \[
    \bigcap_{n=1}^\infty A_n = \{(0, 0)\},
  \]
  but
  \[
    \bigcup_{n=1}^\infty A_n = \left\{(x,y) \vert x^2 + y^2 < 1\right\}.
  \]
\end{example}

\begin{proposition}
  Let \(A\), \(B\) and \(C\) be the random events. Then we have the following properties:
  \begin{enumerate}\index{De Morgan's Laws}
    \item Commutativity: \(A \cup B = B \cup A\), \(A B = B A\)
    \item Associativity: \(A \cup (B \cup C) = (A \cup B) \cup C = A \cup B \cup C\), \(A (B C) = (A B) C = A B C\)
    \item Distributivity: \(A (B \cup C) = A B \cup A C\), \(A \cup (B C) = (A \cup B)(A \cup C)\)
    \item De Morgan's Laws: \((A \cup B)^c = \overline{A} \cap \overline{B}\), \((A \cap B)^c = \overline{A} \cup \overline{B}\), \(\left(\bigcup_{i=1}^n A_i\right)^c = \bigcap_{i=1}^n \overline{A}_i\), \(\left(\bigcap_{i=1}^n A_i\right)^c = \bigcup_{i=1}^n \overline{A}_i\)
    \item \(A \cup B = A \cup \overline{A} B\), \(A \cup B \cup C = A \cup \overline{A} B \cup \overline{A} \overline{B} C\)
  \end{enumerate}
\end{proposition}

\section{Probabilities of Defined on Events}

\begin{definition}[Classical Random Experiment]\index{classical random experiment}\index{simple space}
  A random experiment \(E\) is \emph{classical} if
  \begin{enumerate}
    \item \(E\) contains only different finite sample points, i.e., \(\Omega = \{\omega_1, \omega_2, \dots, \omega_n\}\). We call this kind of sample space \emph{simple space}.
    \item All outcomes are equally likely to occur.
  \end{enumerate}
\end{definition}

If \(A\) is comprised by \(k\) different elementary events, then
\[
  P(A) = \frac{k}{n}.
\]

\begin{definition}[Geometric Random Experiment]\index{geometric random experiment}
  A random experiment \(E\) is \emph{geometric} if
  \begin{enumerate}
    \item The sample space is measurable (such as length, area, volume, etc.) region, i.e., \(0 < \mathrm{L}(\Omega) < \infty\).
    \item The probability of every event \(A \subseteq \Omega\) is proportional to the measure \(L(A)\) and has nothing to do with its position and shape.
  \end{enumerate}
\end{definition}

In a geometric random experiment, the probability of an event \(A\) is given by
\[
  P(A) = \frac{\mathrm{L}(A)}{\mathrm{L}(\Omega)}.
\]

Bertrand's paradox is a famous problem in geometric probability, which shows that the probability of an event can depend on the method of random selection, even if the underlying geometric conditions are the same.

\begin{definition}[Relative Frequency]\index{relative frequency}
  Let \(E\) be an random experiment, \(A\) be an ramdom event. Suppose that \(E\) was repeated \(n\) times under similiar conditions. Let \(f_n(A)\) be the times that \(A\) occurs. The ratio \(\frac{f_n(A)}{n}\) is called the \emph{relative frequency} of event \(A\) in \(n\) trials.
\end{definition}

\begin{theorem}
  For classical random experiment \(E\), the probability has the following properties:
  \begin{enumerate}
    \item for every event \(A\), \(P(A) \geq 0\);
    \item \(P(\Omega) = 1\);
    \item (finite additivity) for every sequence of disjoint events \(A_1, A_2, \dots, A_m\), we have
    \[
      P\left(\bigcup_{i=1}^m A_i\right) = \sum_{i=1}^m P(A_i), \quad m = 1, 2, \dots, n.
    \]
  \end{enumerate}
\end{theorem}

\begin{theorem}
  For geometric random experiment \(E\), the probability has the following properties:
  \begin{enumerate}
    \item for every event \(A\), \(P(A) \geq 0\);
    \item \(P(\Omega) = 1\);
    \item (countable additivity) for every sequence of disjoint events \(A_1, A_2, \dots\), we have
    \[
      P\left(\bigcup_{i=1}^\infty A_i\right) = \sum_{i=1}^\infty P(A_i).
    \]
  \end{enumerate}
\end{theorem}

\begin{theorem}
  For a random experiment \(E\), the frequency has the following properties:
  \begin{enumerate}
    \item for every event \(A\), \(F_n(A) \geq 0\);
    \item \(F_n(\Omega) = 1\);
    \item (finite additivity) for every sequence of disjoint events \(A_1, A_2, \dots, A_m\),
    \[
      F_n\left(\bigcup_{i=1}^m A_i\right) = \sum_{i=1}^m F_n(A_i), \quad m = 1, 2, \dots, n.
    \]
  \end{enumerate}
\end{theorem}

\begin{definition}[\(\sigma\)-algebra]\index{\(\sigma\)-algebra}
  A collection \(\mathcal{F}\) of subsets of \(\Omega\) is a \emph{\(\sigma\)-algebra} if
  \begin{enumerate}
    \item \(\Omega \in \mathcal{F}\);
    \item if \(A \in \mathcal{F}\), then \(\overline{A} \in \mathcal{F}\);
    \item if \(A_n \in \mathcal{F}\) for \(n = 1, 2, \dots\), then \(\bigcup_{n=1}^\infty A_n \in \mathcal{F}\).
  \end{enumerate}
\end{definition}

That is, a \(\sigma\)-algebra is a collection of subsets of \(\Omega\) that contains \(\Omega\) itself, is closed under complementation, and is closed under countable unions.

\begin{example}
  Let the experiment be a coin toss. Then \(\Omega = \{H, T\}\). Let \(\mathcal{F}_1 = \{\varnothing, \{H\}, \{T\}, \Omega\}\) and \(\mathcal{F}_2 = \{\varnothing, \Omega\}\). Then both \(\mathcal{F}_1\) and \(\mathcal{F}_2\) are \(\sigma\)-algebras on \(\Omega\).
\end{example}

\begin{proposition}
  If \(\mathcal{F}\) is a \(\sigma\)-algebra on \(\Omega\), then it has the following properties:
  \begin{enumerate}
    \item \(\varnothing \in \mathcal{F}\);
    \item if \(A, B \in \mathcal{F}\), then \(A B \in \mathcal{F}\) and \(A \setminus B \in \mathcal{F}\);
    \item if \(A_i \in \mathcal{F} (i = 1, 2, \dots, n)\), then \(\bigcup_{i=1}^n A_i \in \mathcal{F}\) and \(\bigcap_{i=1}^n A_i \in \mathcal{F}\);
    \item if \(A_i \in \mathcal{F} (i = 1, 2, \dots)\), then \(\bigcap_{i=1}^\infty A_i \in \mathcal{F}\).
  \end{enumerate}
\end{proposition}

\section{Axiomatic Definition of Probability}

\begin{definition}[Probability]\index{probability}\index{probability space}
  Let \(P(A) (A \in \mathcal{F})\) be a non-negative set function on the \(\sigma\)-algebra \(\mathcal{F}\). Then \(P(A)\) is called the \emph{probability of event \(A\)} if it satisfies the following axioms:
  \begin{enumerate}
    \item \(P(A) \geq 0\) for every event \(A \in \mathcal{F}\);
    \item \(P(\Omega) = 1\);
    \item (countable additivity) for every infinite sequence of countably disjoint events \(A_1, A_2, \dots\), we have
    \[
      P\left(\bigcup_{i=1}^\infty A_i\right) = \sum_{i=1}^\infty P(A_i).
    \]
  \end{enumerate}
  The sets in \(\mathcal{F}\) are called \emph{events}, and \(\mathcal{F}\) is called the \emph{algebra of events}. The triple \((\Omega, \mathcal{F}, P)\) is called a \emph{probability space} or \emph{probability triplet}.
\end{definition}

\begin{proposition}
  Some properties of probability can be derived from the axioms:
  \begin{enumerate}
    \item \(P(\varnothing) = 0\);
    \item for every finite sequence of disjoint events \(A_1, A_2, \dots, A_n\), we have
    \[
      P\left(\bigcup_{i=1}^n A_i\right) = \sum_{i=1}^n P(A_i),
    \]
    \item \(P(A) = 1 - P(\overline{A})\);
    \item if \(A \subseteq B\), then \(P(B \setminus A) = P(B) - P(A)\) and \(P(A) \leq P(B)\);
    \item \(P(A \cup B) = P(A) + P(B) - P(A B)\), similarly, for every finite sequence of events \(A_1, A_2, \dots, A_n\), we have
    \[
      P\left(\bigcup_{i=1}^n A_i\right) = \sum_{i=1}^n P(A_i) - \sum_{1 \leq i < j \leq n} P(A_i A_j) + \dots + (-1)^{n-1} P\left(\bigcap_{i=1}^n A_i\right),
    \]
    \item Continuity: If \(A_n \uparrow A\) (i.e., \(A_1 \subseteq A_2 \subseteq \cdots\) and \(A = \bigcup_{n=1}^\infty A_n\)), then \(P(A) = \lim_{n \to \infty} P(A_n)\); if \(A_n \downarrow A\) (i.e., \(A_1 \supseteq A_2 \supseteq \cdots\) and \(A = \bigcap_{n=1}^\infty A_n\)), then \(P(A) = \lim_{n \to \infty} P(A_n)\).
  \end{enumerate}
\end{proposition}

\begin{proof}
  We only prove the continuity property for \(A_n \uparrow A\), the other case can be proved similarly.
  If \(A_n \uparrow A\), then
  \[
  A = \bigcup_{n=1}^\infty A_n = A_1 \cup A_2 \overline{A}_1 \cup A_3 \overline{A}_2 \cup \cdots,
  \]
  and
  \[
  A_1 \cup A_2 \overline{A}_1 \cup A_3 \overline{A}_2 \cup \dots \cup A_n \overline{A}_{n-1} = A_n.
  \]
  Let \(B_1 = A_1\) and \(B_n = A_n \overline{A}_{n-1}\) for \(n \geq 2\). Then
  \begin{align*}
    P(A) & = P\left(\bigcup_{n=1}^\infty B_n\right) = \sum_{n=1}^\infty P(B_n) = \lim_{n \to \infty} \sum_{i=1}^n P(B_i) \\
         & = \lim_{n \to \infty} P \left(\bigcup_{i=1}^n B_i\right) = \lim_{n \to \infty} P(A_n). \qedhere
  \end{align*}
\end{proof}

\section{Conditional Probability}

\begin{definition}[Conditional Probability]\index{conditional probability}
  Given two events \(A\) and \(B\) with \(P(B) > 0\), the \emph{conditional probability} of \(A\) given \(B\) is defined as
  \[
    P(A \vert B) = \frac{P(A B)}{P(B)}.
  \]
\end{definition}

\begin{theorem}
  If \(P(B) > 0\), then the conditional probability is also a probability, that is,
  \begin{enumerate}
    \item for every event \(A\), \(P(A \vert B) \geq 0\);
    \item \(P(\Omega \vert B) = 1\);
    \item for every infinite sequence of countably disjoint events \(A_1, A_2, \dots\), we have
    \[
      P\left(\bigcup_{i=1}^\infty A_i \vert B\right) = \sum_{i=1}^\infty P(A_i \vert B).
    \]
  \end{enumerate}
\end{theorem}

\begin{proposition}
  Some properties of conditional probability can be derived from the axioms:
  \begin{enumerate}
    \item \(P(\varnothing \vert B) = 0\);
    \item for every finite sequence of disjoint events \(A_1, A_2, \dots, A_n\), we have
    \[
      P\left(\bigcup_{i=1}^n A_i \vert B\right) = \sum_{i=1}^n P(A_i \vert B),
    \]
    \item \(P(A \vert B) = 1 - P(\overline{A} \vert B)\);
    \item if \(A \subseteq C\), then \(P(C \setminus A \vert B) = P(C \vert B) - P(A \vert B)\) and \(P(A \vert B) \leq P(C \vert B)\);
    \item \(P(A \cup C \vert B) = P(A \vert B) + P(C \vert B) - P(A C \vert B)\).
  \end{enumerate}
\end{proposition}

\begin{theorem}
  Assume that \(P(B) > 0\). Then
  \[
    P(A B) = P(B) P(A \vert B).
  \]
\end{theorem}

\begin{theorem}
  Suppose that \(A_1, A_2, \dots, A_n\) are events with \(P\left(\bigcap_{i=1}^{n-1} A_i\right) > 0\). Then
  \[
    P\left(\bigcap_{i=1}^n A_i\right) = P(A_1) P(A_2 \vert A_1) P(A_3 \vert A_1 A_2) \cdots P\left(A_n \vert \bigcap_{i=1}^{n-1} A_i\right).
  \]
\end{theorem}

\begin{example}
  Suppose that an urn box contains $r$ red balls and $b$ black balls ($r, b \geq 2$). We randomly select 1 ball from the box, and then put it back along with $c$ balls of the same color into the box. Then we randomly select another ball from the box. Determine the probability of obtaining the sequence of colors \texttt{red-red-black-black}.
\end{example}

\begin{solution}
  Let \(R_i\) be the event that the \(i\)-th selected ball is red, and let \(B_i\) be the event that the \(i\)-th selected ball is black. Then
  \begin{align*}
    P(R_1 R_2 B_3 B_4) & = P(R_1) P(R_2 \vert R_1) P(B_3 \vert R_1 R_2) P(B_4 \vert R_1 R_2 B_3) \\
                       & = \frac{r}{r + b} \cdot \frac{r + c}{r + b + c} \cdot \frac{b}{r + b + 2c} \cdot \frac{b + c}{r + b + 3c}. \qedhere
  \end{align*}
\end{solution}

\begin{theorem}[Total Probability Theorem]\index{total probability theorem}
  Suppose that \(B_1, B_2, \dots, B_n\) form a partition of the sample space \(\Omega\), and \(P(B_i) > 0\) for \(i = 1, 2, \dots, n\). Then
  \[
    P(A) = \sum_{i=1}^n P(B_i) P(A \vert B_i).
  \]
\end{theorem}

\begin{theorem}[Bayes' Theorem]\index{Bayes' theorem}
  Suppose that \(B_1, B_2, \dots, B_n\) form a partition of the sample space \(\Omega\), and \(P(B_i) > 0\) for \(i = 1, 2, \dots, n\). Let \(A\) be an event with \(P(A) > 0\). Then for every \(i = 1, 2, \dots, n\), we have
  \[
    P(B_i \vert A) = \frac{P(B_i) P(A \vert B_i)}{\sum_{j=1}^n P(B_j) P(A \vert B_j)}.
  \]
\end{theorem}

Bayes' Theorem can be rewritten as
\[
  \underbrace{P(C \vert A)}_{\text{posterior probability}} = \frac{\overbrace{P(C)}^{\text{prior probability}} \overbrace{P(A \vert C)}^{\text{conditional likelihood}}}{\underbrace{P(A)}_{\text{marginal likelihood}}}.
\]

\section{Independence of Events}

\begin{definition}[Independence]\index{independent events}
  Two events \(A\) and \(B\) are \emph{independent} if \(P(A B) = P(A) P(B)\).
\end{definition}

If \(P(B) > 0\), then \(A\) and \(B\) are independent if and only if \(P(A \vert B) = P(A)\). That is, the occurrence of event \(B\) does not affect the probability of event \(A\).

\begin{theorem}
  If \(A\) and \(B\) are independent, then \(A\) and \(\overline{B}\) are independent, \(\overline{A}\) and \(B\) are independent, and \(\overline{A}\) and \(\overline{B}\) are independent.
\end{theorem}

\begin{proof}
  We only prove that \(A\) and \(\overline{B}\) are independent, the other cases can be proved similarly. Since \(A\) and \(B\) are independent, we have
  \[
    P(A \overline{B}) = P(A) - P(A B) = P(A) - P(A) P(B) = P(A)(1 - P(B)) = P(A) P(\overline{B}). \qedhere
  \]
\end{proof}

\begin{example}
  Suppose that \(P(A) = 1\) and \(B \in \mathcal{F}\), we want to determine whether \(A\) and \(B\) are independent. Since \(P(A) = 1\), we have
  \[
    0 \leq P(\overline{A} B) \leq P(\overline{A}) = 0,
  \]
  which implies \(P(\overline{A} B) = 0\).
  Therefore, \(P(A B) = P(B) - P(\overline{A} B) = P(B) = P(A) P(B)\), and \(A\) and \(B\) are independent.
\end{example}

\begin{definition}[Mutual Independence]\index{mutual independence}\index{pairwise independence}
  Three events \(A\), \(B\) and \(C\) are \emph{mutually independent} (or \emph{independent}) if
  \begin{align*}
    P(A B) & = P(A) P(B), \\
    P(A C) & = P(A) P(C), \\
    P(B C) & = P(B) P(C), \\
    P(A B C) & = P(A) P(B) P(C).
  \end{align*}
  Events \(A\), \(B\) and \(C\) are \emph{pairwise independent} if only the first three conditions hold.
\end{definition}

\begin{definition}[Mutual Independence]\index{mutual independence}
  Events \(A_1, A_2, \dots, A_n\) are \emph{independent} if for every subset \(\{i_1, i_2, \dots, i_k\} \subseteq \{1, 2, \dots, n\}\) with \(k \geq 2\), we have
  \[
    P\left(\bigcap_{j=1}^k A_{i_j}\right) = \prod_{j=1}^k P(A_{i_j}).
  \]
\end{definition}

\begin{theorem}
  Events \(A_1, A_2, \dots, A_n\) are independent if and only if for every subset \(\{i_1, i_2, \dots, i_k\} \subseteq \{1, 2, \dots, n\}\) with \(k \geq 2\), we have
  \[
    P\left(\bigcap_{j=1}^k A_{i_j}^*\right) = \prod_{j=1}^k P(A_{i_j}^*),
  \]
  where \(A_i^*\) denotes \(A_i\) or \(\overline{A}_i\) for \(i = 1, 2, \dots, n\).
\end{theorem}

If events \(A_1, A_2, \dots, A_n\) are independent, by De Morgan's Laws, we have
\[
  P\left(\bigcup_{i=1}^n A_i\right) = 1 - P\left(\bigcap_{i=1}^n \overline{A}_i\right) = 1 - \prod_{i=1}^n (1 - P(A_i)),
\]
and
\[
  P\left(\bigcup_{i=1}^n \overline{A}_i\right) = 1 - P\left(\bigcap_{i=1}^n A_i\right) = 1 - \prod_{i=1}^n P(A_i).
\]

\begin{definition}[Bernoulli Trial]\index{Bernoulli trial}\index{Bernoulli process}
  A Bernoulli trial is a random experiment with only two possible outcomes, usually denoted as "success" and "failure". The probability of success is denoted by \(p\) and the probability of failure is denoted by \(q = 1 - p\). A sequence of independent Bernoulli trials is called a Bernoulli process.
\end{definition}

\begin{theorem}
  The probability that the outcome of an experiment that consists of \(n\) Bernoulli trials contains exactly \(k\) successes is given by the binomial distribution:
  \[
    P(X = k) = \binom{n}{k} p^k q^{n-k}, \quad k = 0, 1, \dots, n.
  \]
\end{theorem}